
\learninguniverse{
title={AI & Machine Learning Mastery}
description={Comprehensive AI & ML course covering everything from foundations to deep learning!}
difficulty={Beginner to Advanced}
estimatedHours={80}
skills={Python,Machine Learning,Deep Learning,NumPy,Pandas,Matplotlib,TensorFlow,PyTorch}
}

\track{
title={Machine Learning Foundations}
description={Build your core AI & ML knowledge!}
learningoutcomes={Understand AI fundamentals, master ML basics}

\module{
title={Introduction to AI}
description={Start your AI journey here!}
estimatedHours={10}

\lesson{
title={What is AI}
overviewmarkdown={
# What is Artificial Intelligence

Artificial Intelligence (AI) is the simulation of human intelligence in machines that are programmed to think like humans and mimic their actions.
}
}

\lesson{
title={Types of AI}
}

\lesson{
title={Real World Applications}
}

\video{
type={youtube}
url={https://www.youtube.com/watch?v=J4Q39xW3V8Y}
}
\video{
type={youtube}
url={https://www.youtube.com/watch?v=O9v47V0B514}
}
\video{
type={youtube}
url={https://www.youtube.com/watch?v=W6Z5Z5J5Z5J}
}

\practice{
language={python}
startercode={
print("Hello AI 1")
}
expectedoutput={Hello AI 1}
}
\practice{
language={python}
startercode={
print("Hello AI 2")
}
expectedoutput={Hello AI 2}
}
\practice{
language={python}
startercode={
print("Hello AI 3")
}
expectedoutput={Hello AI 3}
}

\quiz{
question={What is Artificial Intelligence?}
optionA={Type of robot}
optionB={Simulation of human intelligence in machines}
optionC={Programming language}
optionD={Database type}
correct={B}
}
\quiz{
question={Which is a type of AI?}
optionA={Narrow AI}
optionB={Wide AI}
optionC={Long AI}
optionD={Tall AI}
correct={A}
}
\quiz{
question={What does NLP stand for?}
optionA={Natural Language Processing}
optionB={Network Language Protocol}
optionC={Neural Learning Program}
optionD={New Learning Platform}
correct={A}
}

\project{
title={Chatbot}
colab={https://colab.research.google.com/drive/1example}
instructions={Build a chatbot}
}
\project{
title={Image Recognition}
colab={https://colab.research.google.com/drive/2example}
instructions={Build an image recognizer}
}
\project{
title={Sentiment Analysis}
colab={https://colab.research.google.com/drive/3example}
instructions={Build a sentiment analyzer}
}

\resource{
type={website}
title={NVIDIA AI Blog}
url={https://blogs.nvidia.com}
}
\resource{
type={github}
title={Awesome AI}
url={https://github.com/owainlewis/awesome-artificial-intelligence}
}
\resource{
type={pdf}
title={AI Foundations}
url={https://example.com/ai-foundations.pdf}
}
}

\module{
title={Machine Learning Basics}
description={Learn core ML concepts!}
estimatedHours={12}

\lesson{
title={What is ML}
}
\lesson{
title={Supervised Learning}
}
\lesson{
title={Unsupervised Learning}
}

\video{
type={youtube}
url={https://www.youtube.com/watch?v=ukzFI9rgwfU}
}
\video{
type={youtube}
url={https://www.youtube.com/watch?v=Gv9_4yMHFhI}
}
\video{
type={youtube}
url={https://www.youtube.com/watch?v=nKxmqSCGmek}
}

\practice{
language={python}
startercode={
import numpy as np
from sklearn.linear_model import LinearRegression
X = np.array([[1], [2], [3]])
y = np.array([2, 4, 6])
model = LinearRegression().fit(X, y)
print("Predict 4:", model.predict([[4]]))
}
expectedoutput={Predict 4: [8.]}
}
\practice{
language={python}
startercode={
import numpy as np
from sklearn.svm import SVC
X = np.array([[1, 2], [3, 4], [5,6]])
y = np.array([0,0,1])
model = SVC().fit(X, y)
print("Predict [1,1]:", model.predict([[1,1]]))
}
expectedoutput={Predict [1,1]: [0]}
}
\practice{
language={python}
startercode={
import numpy as np
from sklearn.cluster import KMeans
X = np.array([[1, 2], [3,4], [100,101], [102,103]])
kmeans = KMeans(n_clusters=2, random_state=0).fit(X)
print(kmeans.labels_)
}
expectedoutput={[0 0 1 1]}
}

\quiz{
question={What is Machine Learning?}
optionA={Programming with explicit instructions}
optionB={Learning from data}
optionC={Type of computer}
optionD={Type of internet}
correct={B}
}
\quiz{
question={Which uses labeled data?}
optionA={Unsupervised}
optionB={Supervised}
optionC={Reinforcement}
optionD={None}
correct={B}
}
\quiz{
question={Clustering uses?}
optionA={Supervised}
optionB={Unsupervised}
optionC={Reinforcement}
optionD={All}
correct={B}
}

\project{
title={Regression}
colab={https://colab.research.google.com/drive/4example}
instructions={Build regression model}
}
\project{
title={Classification}
colab={https://colab.research.google.com/drive/5example}
instructions={Build classifier}
}
\project{
title={Clustering}
colab={https://colab.research.google.com/drive/6example}
instructions={Build cluster model}
}

\resource{
type={github}
title={Awesome ML}
url={https://github.com/josephmisiti/awesome-machine-learning}
}
\resource{
type={website}
title={ML Mastery}
url={https://machinelearningmastery.com}
}
\resource{
type={pdf}
title={ML Textbook}
url={https://example.com/ml-textbook.pdf}
}
}
}

\track{
title={Data Science Essentials}
description={Learn data science tools and preprocessing!}

\module{
title={Python for ML}
description={Master NumPy, Pandas, and Matplotlib!}
estimatedHours={15}

\lesson{
title={NumPy}
}
\lesson{
title={Pandas}
}
\lesson{
title={Matplotlib}
}

\video{
type={youtube}
url={https://www.youtube.com/watch?v=8MfGp7Fm8}
}
\video{
type={youtube}
url={https://www.youtube.com/watch?v=zgbHp65XnN4}
}
\video{
type={youtube}
url={https://www.youtube.com/watch?v=3ZZzL5K-8}
}

\practice{
language={python}
startercode={
import numpy as np
arr = np.arange(10)
print("NumPy:", arr)
}
expectedoutput={NumPy: [0 1 2 3 4 5 6 7 8 9]}
}
\practice{
language={python}
startercode={
import pandas as pd
df = pd.DataFrame({"x": [1,2,3], "y": [4,5,6]})
print(df)
}
expectedoutput={   x  y
0  1  4
1  2  5
2  3  6}
}
\practice{
language={python}
startercode={
import matplotlib.pyplot as plt
plt.plot([1,2,3], [1,4,9])
plt.savefig("plot.png")
print("Plot saved")
}
expectedoutput={Plot saved}
}

\quiz{
question={NumPy is used for?}
optionA={Web dev}
optionB={Numerical computation}
optionC={Game dev}
optionD={Design}
correct={B}
}
\quiz{
question={Pandas is for?}
optionA={Web scraping}
optionB={Data frames}
optionC={ML}
optionD={Images}
correct={B}
}
\quiz{
question={Matplotlib plots?}
optionA={Data}
optionB={Videos}
optionC={Sounds}
optionD={Files}
correct={A}
}

\project{
title={NumPy Exercises}
colab={https://colab.research.google.com/drive/7example}
instructions={Solve NumPy problems}
}
\project{
title={Pandas Exercises}
colab={https://colab.research.google.com/drive/8example}
instructions={Solve Pandas problems}
}
\project{
title={Data Visualization}
colab={https://colab.research.google.com/drive/9example}
instructions={Create visualizations}
}

\resource{
type={github}
title={NumPy Docs}
url={https://github.com/numpy/numpy}
}
\resource{
type={github}
title={Pandas Docs}
url={https://github.com/pandas-dev/pandas}
}
\resource{
type={website}
title={Matplotlib Docs}
url={https://matplotlib.org}
}
}

\module{
title={Data Preprocessing}
description={Clean and prepare your data!}
estimatedHours={10}

\lesson{
title={Missing Values}
}
\lesson{
title={Feature Scaling}
}
\lesson{
title={Encoding}
}

\video{
type={youtube}
url={https://www.youtube.com/watch?v=k932K15vI}
}
\video{
type={youtube}
url={https://www.youtube.com/watch?v=q915v1o5I}
}
\video{
type={youtube}
url={https://www.youtube.com/watch?v=k932K15vI_2}
}

\practice{
language={python}
startercode={
import pandas as pd
import numpy as np
df = pd.DataFrame({"x": [1, np.nan,3]})
df = df.fillna(0)
print(df)
}
expectedoutput={     x
0  1.0
1  0.0
2  3.0}
}
\practice{
language={python}
startercode={
from sklearn.preprocessing import StandardScaler
import numpy as np
data = np.array([[1,2],[3,4]])
scaler = StandardScaler()
print(scaler.fit_transform(data))
}
expectedoutput={[[-1. -1.]
[ 1.  1.]]}
}
\practice{
language={python}
startercode={
from sklearn.preprocessing import LabelEncoder
import pandas as pd
data = pd.Series(["cat","dog","cat"])
le = LabelEncoder()
print(le.fit_transform(data))
}
expectedoutput={[0 1 0]}
}

\quiz{
question={What is preprocessing?}
optionA={Cleaning data}
optionB={Training model}
optionC={Testing model}
optionD={Deploying model}
correct={A}
}
\quiz{
question={Scaling helps?}
optionA={Model convergence}
optionB={Model speed}
optionC={Model size}
optionD={Nothing}
correct={A}
}
\quiz{
question={LabelEncoder for?}
optionA={Text to numbers}
optionB={Numbers to text}
optionC={Images}
optionD={Sounds}
correct={A}
}

\project{
title={Preprocessing}
colab={https://colab.research.google.com/drive/Aexample}
instructions={Clean dataset}
}
\project{
title={Feature Engineering}
colab={https://colab.research.google.com/drive/Bexample}
instructions={Create features}
}
\project{
title={Data Pipeline}
colab={https://colab.research.google.com/drive/Cexample}
instructions={Build pipeline}
}

\resource{
type={github}
title={Feature Engineering}
url={https://github.com/FeatureEngineering/FeatureEngineering}
}
\resource{
type={website}
title={Data Preprocessing Guide}
url={https://www.kaggle.com/code/rtatman/data-cleaning-challenge-handling-missing-values}
}
\resource{
type={pdf}
title={Preprocessing Textbook}
url={https://example.com/preprocessing-book.pdf}
}
}
}

\track{
title={Machine Learning Algorithms}
description={Dive into regression and classification!}

\module{
title={Regression}
description={Learn linear and multiple regression!}
estimatedHours={12}

\lesson{
title={Linear Regression}
}
\lesson{
title={Multiple Regression}
}

\video{
type={youtube}
url={https://www.youtube.com/watch?v=zPG1n1B0Y}
}
\video{
type={youtube}
url={https://www.youtube.com/watch?v=nk2CQITm_eA}
}

\practice{
language={python}
startercode={
import numpy as np
from sklearn.linear_model import LinearRegression
X = np.array([[1,2],[3,4],[5,6]])
y = np.array([3,7,11])
model = LinearRegression()
model.fit(X,y)
print(model.predict([[7,8]]))
}
expectedoutput={[15.]}
}
\practice{
language={python}
startercode={
import numpy as np
from sklearn.datasets import fetch_california_housing
X,y = fetch_california_housing(return_X_y=True)
from sklearn.linear_model import LinearRegression
model = LinearRegression()
model.fit(X,y)
print(model.score(X,y))
}
expectedoutput={0.606232928751043}
}

\quiz{
question={Regression predicts?}
optionA={Continuous values}
optionB={Categories}
optionC={Both}
optionD={None}
correct={A}
}
\quiz{
question={Multiple regression uses?}
optionA={One feature}
optionB={Multiple features}
optionC={No features}
optionD={Images}
correct={B}
}

\project{
title={House Price Prediction}
colab={https://colab.research.google.com/drive/Dexample}
instructions={Predict house prices}
}
\project{
title={Stock Price Prediction}
colab={https://colab.research.google.com/drive/Eexample}
instructions={Predict stock prices}
}

\resource{
type={github}
title={Regression Resources}
url={https://github.com/topics/regression}
}
\resource{
type={website}
title={Regression Tutorial}
url={https://towardsdatascience.com/regression-analysis-explained}
}
}

\module{
title={Classification}
description={Learn logistic regression and decision trees!}
estimatedHours={14}

\lesson{
title={Logistic Regression}
}
\lesson{
title={Decision Trees}
}

\video{
type={youtube}
url={https://www.youtube.com/watch?v=yIYKR4sgzI8}
}
\video{
type={youtube}
url={https://www.youtube.com/watch?v=7VeUPuFGJHk}
}

\practice{
language={python}
startercode={
import numpy as np
from sklearn.linear_model import LogisticRegression
X = np.array([[1,2],[3,4],[5,6],[7,8]])
y = np.array([0,0,1,1])
model = LogisticRegression()
model.fit(X,y)
print(model.predict([[1,1]]))
}
expectedoutput={[0]}
}
\practice{
language={python}
startercode={
import numpy as np
from sklearn.tree import DecisionTreeClassifier
X = np.array([[1,2],[3,4],[5,6],[7,8]])
y = np.array([0,0,1,1])
model = DecisionTreeClassifier()
model.fit(X,y)
print(model.predict([[1,1]]))
}
expectedoutput={[0]}
}

\quiz{
question={Logistic regression for?}
optionA={Classification}
optionB={Regression}
optionC={Clustering}
optionD={None}
correct={A}
}
\quiz{
question={Decision trees split data using?}
optionA={Features}
optionB={Randomly}
optionC={Colors}
optionD={Images}
correct={A}
}

\project{
title={Iris Classification}
colab={https://colab.research.google.com/drive/Fexample}
instructions={Classify Iris flowers}
}
\project{
title={Breast Cancer Detection}
colab={https://colab.research.google.com/drive/Gexample}
instructions={Detect cancer}
}

\resource{
type={github}
title={Classification}
url={https://github.com/topics/classification}
}
\resource{
type={website}
title={Classification Tutorial}
url={https://towardsdatascience.com/classification-in-machine-learning}
}
}
}

\track{
title={Deep Learning}
description={Enter the world of neural networks!}

\module{
title={Neural Networks}
description={Learn core neural network concepts!}
estimatedHours={12}

\lesson{
title={Perceptrons}
}
\lesson{
title={Activation Functions}
}
\lesson{
title={Forward Propagation}
}

\video{
type={youtube}
url={https://www.youtube.com/watch?v=aircA707Y}
}
\video{
type={youtube}
url={https://www.youtube.com/watch?v=umsmS2l5J}
}
\video{
type={youtube}
url={https://www.youtube.com/watch?v=O5xqq058Mxo}
}

\practice{
language={python}
startercode={
import numpy as np
def sigmoid(x): return 1/(1+np.exp(-x))
print(sigmoid(0))
}
expectedoutput={0.5}
}
\practice{
language={python}
startercode={
import numpy as np
def relu(x): return max(0,x)
print(relu(-5), relu(5))
}
expectedoutput={0 5}
}
\practice{
language={python}
startercode={
import numpy as np
def sigmoid(x): return 1/(1+np.exp(-x))
W = np.array([0.5, 0.5])
b = 0
X = np.array([1,1])
print(sigmoid(W.dot(X)+b))
}
expectedoutput={0.7310585786300049}
}

\quiz{
question={Perceptron basics of?}
optionA={Input, weights, bias}
optionB={Images}
optionC={Videos}
optionD={Files}
correct={A}
}
\quiz{
question={Activation functions add?}
optionA={Nonlinearity}
optionB={Linearity}
optionC={Noise}
optionD={Nothing}
correct={A}
}
\quiz{
question={Forward propagation computes?}
optionA={Output from input}
optionB={Input from output}
optionC={Random numbers}
optionD={Images}
correct={A}
}

\project{
title={Neural Network from Scratch}
colab={https://colab.research.google.com/drive/Hexample}
instructions={Build neural network}
}
\project{
title={Activation Functions}
colab={https://colab.research.google.com/drive/Iexample}
instructions={Implement activation functions}
}
\project{
title={XOR Problem}
colab={https://colab.research.google.com/drive/Jexample}
instructions={Solve XOR problem}
}

\resource{
type={github}
title={Neural Networks}
url={https://github.com/topics/neural-network}
}
\resource{
type={website}
title={Deep Learning Book}
url={https://www.deeplearningbook.org}
}
\resource{
type={pdf}
title={Neural Networks and Deep Learning}
url={http://neuralnetworksanddeeplearning.com/}
}
}

\module{
title={TensorFlow & PyTorch}
description={Build and train models with TensorFlow and PyTorch!}
estimatedHours={15}

\lesson{
title={Building Networks}
}
\lesson{
title={Training Models}
}

\video{
type={youtube}
url={https://www.youtube.com/watch?v=tPYj3fFJGjk}
}
\video{
type={youtube}
url={https://www.youtube.com/watch?v=IC0_FRiXwKo}
}

\practice{
language={python}
startercode={
import tensorflow as tf
model = tf.keras.Sequential([
    tf.keras.layers.Dense(32, activation='relu', input_shape=(784,)),
    tf.keras.layers.Dense(10, activation='softmax')
])
model.compile(optimizer='adam', loss='categorical_crossentropy', metrics=['accuracy'])
print("TF model built")
}
expectedoutput={TF model built}
}
\practice{
language={python}
startercode={
import torch
import torch.nn as nn
model = nn.Sequential(
    nn.Linear(784, 32),
    nn.ReLU(),
    nn.Linear(32,10)
)
print("PyTorch model built")
}
expectedoutput={PyTorch model built}
}

\quiz{
question={TensorFlow made by?}
optionA={Google}
optionB={Facebook}
optionC={OpenAI}
optionD={Microsoft}
correct={A}
}
\quiz{
question={PyTorch made by?}
optionA={Facebook}
optionB={Google}
optionC={OpenAI}
optionD={Microsoft}
correct={A}
}

\project{
title={TensorFlow MNIST}
colab={https://colab.research.google.com/drive/Kexample}
instructions={Train TensorFlow model on MNIST}
}
\project{
title={PyTorch MNIST}
colab={https://colab.research.google.com/drive/Lexample}
instructions={Train PyTorch model on MNIST}
}

\resource{
type={github}
title={TensorFlow}
url={https://github.com/tensorflow/tensorflow}
}
\resource{
type={github}
title={PyTorch}
url={https://github.com/pytorch/pytorch}
}
\resource{
type={website}
title={TensorFlow Docs}
url={https://www.tensorflow.org}
}
}
}
